\section{\textsc{MobileGym++} Playground}
\label{sec:mobilegym++}
\label{sec:method}

\subsection{A gap in controllable commercial applications}
\label{sec:mgpp-env}

Mobile agents are increasingly moving toward real-world workflows across commercial applications, such as travel booking on Trip and online shopping on Taobao, where tasks are often long-horizon and grounded in real user needs. However, training and evaluation environments for such commercial applications remain largely missing, leaving agents with limited opportunities to practice and be evaluated in realistic mobile-use scenarios.

Real-device or Android-emulator setups introduce login friction, irreversible side effects, and costly reset procedures that make large-scale training impractical~\citep{rawles2025androidworld,kong2026mobileworld}. Prior simulated-app work, meanwhile, covers only a narrow range of commercial domains~\citep{wu2026mobilegym,trivedi2024appworld}.

We built MobileGym++, a browser-hosted mobile environment that hosts realistic commercial workflows in controllable simulated applications. The runtime extends MobileGym~\citep{wu2026mobilegym}, preserving its advantages in reset and programmatic evaluation while substantially expanding domain coverage, the number and richness of applications, and support for realistic cross-app workflows. Observations consist of raw screen pixels; there is no accessibility tree and the agent cannot read business state directly. Device state is fully serializable, enabling deterministic reset and parallel instantiation. We newly built 35 commercial applications in this round, bringing the total to 63 installable clients spanning travel, shopping, finance, content, social, office, and local services. Each application maintains persistent business state---orders, budgets, itineraries---that evolves through interaction and supports programmatic task verification.

% Figure~\ref{fig:mgpp-overview} was split out of this section; the overview
% panel now lives in the introduction. Keep the float commented until a
% replacement is dropped into figures/.

% \begin{table*}[t]
% \caption{Comparison of mobile-agent benchmarks on scale, interaction, and evaluation. MobileGym++ uses a 63-application simulated environment of commercial-style apps. The Platform column distinguishes API, Android, and simulated settings. Following recent benchmark comparisons, capabilities are shown as explicit columns. $\checkmark$ denotes supported, $\times$ unsupported, and --- not reported or not applicable. The two MobileGym rows share the same reproducible simulation and deterministic state-based evaluation; scores are omitted because the platforms and task definitions do not share a protocol.}
% \label{tab:mgpp-stack-compare}
% \centering
% \scriptsize
% \setlength{\tabcolsep}{3.0pt}
% \renewcommand{\arraystretch}{1.16}
% \resizebox{\linewidth}{!}{%
% \begin{tabular}{@{}l l r r c c c c c@{}}
% \toprule
%  & & \multicolumn{2}{c}{Scale} & \multicolumn{3}{c}{Interaction} & \multicolumn{2}{c}{Evaluation} \\
% \cmidrule(lr){3-4} \cmidrule(lr){5-7} \cmidrule(lr){8-9}
% Benchmark & Platform & \# Apps & \# Tasks & Cross-app & GUI & Native tools / MCP & Reproducible & Deterministic eval. \\
% \midrule
% AppWorld~\citep{trivedi2024appworld} & API & 9 & 750 / 250 sc. & $\checkmark$ & $\times$ & --- & $\checkmark$ & $\checkmark$ \\
% AndroidLab~\citep{xu2025androidlab} & Android & 9 & 138 & --- & $\checkmark$ & --- & --- & $\checkmark$ \\
% AndroidWorld~\citep{rawles2025androidworld} & Android & 20 & 116 tmpl. & $\checkmark$ & $\checkmark$ & --- & $\checkmark$ & $\checkmark$ \\
% SPA-Bench~\citep{chen2024spa} & Android & 66 & 340 & $\checkmark$ & $\checkmark$ & --- & $\checkmark$ & $\times$ \\
% MobileWorld~\citep{kong2026mobileworld} & Android & 20 & 201 & $\checkmark$ & $\checkmark$ & $\checkmark$ & $\checkmark$ & $\checkmark$ \\
% MobileGym-Bench~\citep{wu2026mobilegym} & Sim. & 28 & 416 tmpl. & $\checkmark$ & $\checkmark$ & --- & $\checkmark$ & $\checkmark$ \\
% \textbf{MobileGym++ Bench (ours)} & \textbf{Sim.} & \textbf{63$^\dagger$} & \textbf{215} & \textbf{$\checkmark$} & \textbf{$\checkmark$} & \textbf{$\checkmark$} & \textbf{$\checkmark$} & \textbf{$\checkmark$} \\
% \bottomrule
% \end{tabular}%
% }
% \vspace{2pt}
% \begin{flushleft}
% \footnotesize
% $^\dagger$ The environment contains 63 installable applications. --- means that the cited benchmark does not report the capability in a comparable form. Deterministic evaluation refers to an automatic, reproducible success check, not a model judge.
% \end{flushleft}
% \end{table*}

\begin{table*}[t]
\caption{Comparison of mobile agent environments across scale, interaction, and controllability. \ding{51} denotes supported, \ding{55} unsupported, and $\sim$ indicates limited system level cross-app interaction support in MobileGym. \# Comm. Apps and \# Tasks denote the number of covered commercial apps and benchmark tasks, respectively. MobileGym++ provides diverse, high-fidelity commercial-style applications with full controllability, while introducing an experimental app-native tool interface.}
\label{tab:mgpp-stack-compare}
\centering
\scriptsize
\setlength{\tabcolsep}{3.0pt}
\renewcommand{\arraystretch}{1.16}
\resizebox{\linewidth}{!}{%
\begin{tabular}{@{}l l c c c c c c@{}}
\toprule
\multirow{2}{*}[-0.6ex]{Environment} 
& \multirow{2}{*}[-0.6ex]{Platform}
& \multicolumn{2}{c}{Scale} 
& \multicolumn{3}{c}{Interaction} 
& \multirow{2}{*}[-0.6ex]{\shortstack{Full\\Controllability}} \\
\cmidrule(lr){3-4} 
\cmidrule(lr){5-7}
& & \# Comm. Apps & \# Tasks  
& Cross-app & GUI & App-native Tools 
& \\
\midrule
AppWorld~\citep{trivedi2024appworld} & API & 9 & 750  
& \ding{51} & \ding{55} & \ding{51} & \ding{55} \\

AndroidLab~\citep{xu2025androidlab} & Emulator & 5 & 138  
& \ding{55} & \ding{51} & \ding{55} & \ding{55} \\

AndroidWorld~\citep{rawles2025androidworld} & Emulator & 0 & 116  
& \ding{51} & \ding{51} & \ding{55} & \ding{55} \\

SPA-Bench~\citep{chen2024spa} & Real-device & 58 & 340  
& \ding{51} & \ding{51} & \ding{55} & \ding{55} \\

MobileWorld~\citep{kong2026mobileworld} & Emulator & 4 & 201  
& \ding{51} & \ding{51} & \ding{55} & \ding{55} \\

MobileGym~\citep{wu2026mobilegym} & Simulator & 12 & 416  
& $\sim$ & \ding{51} & \ding{55} & \ding{51} \\

\textbf{MobileGym++ (ours)} & Simulator & \textbf{35} & \textbf{215}  
& \greencheck & \greencheck & \greencheck & \greencheck \\

\bottomrule
\end{tabular}%
}
\vspace{2pt}
\end{table*}

\subsection{35 high-fidelity commercial applications}
\label{sec:mgpp-apps}

The 35 applications are structurally faithful replicas of their real-world counterparts. Construction starts from the source application's data schema and interaction graph, specified before any interface is rendered, so business-state richness and navigation depth are architectural properties of each replica. Human annotators verify each replica against the live device; the implementation is corrected until the specified routes, state changes, and checks pass (Appendix~\ref{app:mgpp-build}).

Across these clients, the median application exposes 37 routes, 165.5 interactive actions, and 87 page transitions, with a median interaction depth of 3 from the home screen. Relative to the 12 prior MobileGym business apps retained in the same environment, the newly built clients reach $1.48\times$ the interactive actions and $1.91\times$ the page transitions (Figure~\ref{fig:mgpp-app-scale}). The tool channel is described in Section~\ref{sec:mgpp-tools}.

\begin{figure}[t]
\centering
\includegraphics[width=\linewidth]{figures/fig_app_richness_act_trans.pdf}
\caption{App-level interaction richness. Each marker is one app, plotted by its interactive actions and declared page transitions; marker area encodes the number of Bench tasks referencing the app. Filled circles are the newly built commercial clients (colored by domain); open circles are 12 prior MobileGym business apps. Crosses mark group medians: the new clients reach $1.48\times$ the actions and $1.91\times$ the transitions of the prior ones.}
\label{fig:mgpp-app-scale}
\end{figure}


\subsection{Realistic cross-application workflows}
\label{sec:mgpp-crossapp}

Single-application benchmarks cannot detect whether an agent can identify a required value from an external source, carry it across an application boundary, and write it into the correct field---a routine real-device workflow that demands both information retrieval and persistence across isolated state stores.

MobileGym++ places commercial applications on a single resettable device whose shared state is fully serializable, so the exact multi-app starting configuration can be restored for every training episode or evaluation run. Each application maintains its own live business objects---orders, expenses, itineraries, contacts---which evolve through on-screen actions. Clipboard, system broadcasts, and inter-app navigation connect the applications, so a workflow can leave one client and write into another's store rather than merely forwarding a message.

\subsection{Hybrid GUI and app-native tool use}
\label{sec:mgpp-tools}

Commercial applications commonly expose structured lookup and write interfaces alongside their graphical screens, reflecting that the corresponding records would otherwise be reached only through deep GUI navigation. Hybrid interaction---mixing on-screen actions with application-level calls---already appears in prior work~\citep{kong2026mobileworld}. We expose app-native tools according to a simple, principled design: each tool covers structured search or structured writes over an application's own business objects, and every tool call updates the same persistent store as the corresponding GUI action. This equivalence is what enables a controlled comparison---the GUI-only and Hybrid settings share identical tasks, initial states, and terminal checks, so any performance difference isolates the contribution of tool access rather than task selection.

52 applications expose such tools in the foreground, occupying the same action slot as a tap, swipe, or text input. GUI-only remains a complete setting: every task in the benchmark can be completed through the graphical interface. Hybrid adds a second execution path over the same tasks, initial states, and terminal checks, without changing what counts as success.

\subsection{\textsc{MobileGym++ Bench}}
\label{sec:mgpp-bench}

MobileGym++ Bench evaluates whether an agent can complete a commercial write inside resettable simulated applications, when the required value must be recovered through deep in-app navigation and, in cross-app tasks, from another application's live state rather than from the instruction. Existing mobile-agent benchmarks already evaluate single-app GUI operations, and several include cross-application workflows~\citep{rawles2025androidworld,kong2026mobileworld}, but those settings typically terminate on device or backend state rather than on the application's own business objects.

\paragraph{Task Design.}
The primary source of difficulty is commercial navigation and information flow. A single-app task requires the agent to locate, filter, and fill fields across a deep page stack before a business object is written. A cross-application task further requires a live value from another application's store to be identified, carried across the application boundary, and written into the correct field. Instructions are phrased as user goals and do not name widgets, tools, or checker fields; the terminal is a write in the application's own business state.

\paragraph{Scale and Composition.}
The benchmark comprises 215 tasks spanning seven business domains: 109 tasks involve a single application, 74 involve two, 27 involve three, and 5 involve four. Tasks are distributed across applications and business objects so that aggregate performance is not dominated by a small subset of clients.

\paragraph{Evaluation Protocol.}
Each task specifies a user goal, a bound initial state, and a programmatic checker over the before/after diff of application business state. The checker does not inspect the trajectory or screenshots and does not rely on a model judge. GUI-only and Hybrid are evaluated on the same tasks, initial states, and checks. The primary metric is success rate.
